\scanpage{108}
% archive foliation: 107
de \deleted{niveau} poids $\le2q$, de niveau $\le q$. Donc il en est de même de $E^{\unsure{rs|pq}}_2$ si $p+q=i$, $p\ne0$, \inserted{et on trouve des \deleted{\unsure{poids}} poids $\le2(i-1)<2i$, des niveaux $\le i-1<i$}. Reste à examiner le terme $E^{0,i}_2\subset H^i(X')$; comme $X'$ est lisse, il est donc effectif, de poids $\le2i$, de niveau $\le i$. On a gagné ! Si on note que si $i>n$, on a $H^i(X)=0$ ($X$ affine !).

\deleted{\unsure{À priori}, $H^i(X)$ est de niveau $\le n=\dim X$. \unsure{Le tout} \illegible{4} $E^{pq}_2$ \unsure{est effectif}. \unsure{filtration} \illegible{2} \unsure{niveau}.}

\textcircled{B} \uline{Coefficients constants tordus sur $X$ lisse}, \inserted{séparé} $F$ \inserted{effectif} de poids et niveau $j$.

\textcircled{1} On suppose que les coefficients envisagés proviennent d'un motif lisse sur $X$, et \inserted{par dévissage (possible ici ?)} \unsure{on ne regarde le} cas où ce dernier est un facteur direct d'un $R^jf_*(Q_{l Y})$, avec $f:\deleted{\unsure{Z}}\mathbin{-}X$ morphisme projectif et lisse. \inserted{On peut si on veut supposer $f$ de dim relative $j$.} On est donc \unsure{ramené} au cas où
\[
F=R^jf_*(Q_{l Y}).
\]
Considérons la suite spectrale
\[
H^{\unsure{*}}(Y)\deleted{\unsure{a}}\Longleftarrow
\inserted{E^{pq}_2=}
H^p(X,R^qf_*(Y))
\]
\margin{$\unsure{p}$}

\uuline{D'après Deligne}, cette suite spectrale dégénère. Donc $H^i(X,R^jf_*(Q_l(Y)))$ \inserted{et par suite $H^i(X,F)$} apparait comme un \uline{morceau} de $H^{i+j}(Y)$. On trouve que l'on a un \deleted{\unsure{module}} \uline{effectif}, de poids \deleted{\unsure{\ne}} compris entre
$\deleted{2(n-j)\text{ et }2(i+j)}$
\uline{$i+j$} et $2(n+j)=\uline{2n+2j}$ (où $n=\dim X$);
les niveaux sont resp. $\le(2n-i+j),\ldots,0$.
Si $i\le n$, il n'y a de poids que \unsure{ceux énumérés} $i+j,i+j+1,\ldots,(2(i+j))$, de \deleted{poids} \inserted{niveaux} respectifs
$\deleted{0,\ldots,i+j}$ \inserted{$i+j,i+j-1,\ldots,0$}.

\margin{\fbox{\parbox{.27\linewidth}{peut-on remplacer $2n+2j$ par $2n+j$ ?}}}
\margin{\fbox{\parbox{.27\linewidth}{peut-on remplacer $2i+2j$ par $2i+j$ ? \deleted{\unsure{Cela fait à croire}}}}}

Résultats duaux pour $H^i_!(F)$, il est effectif à poids $\le i+j$, et \deleted{à niveaux} si $i\ge n$, il est à niveaux croiss. de $0$ à $\deleted{2n-i+j}2(n+j)-(i+j)=2n-i+j$, pour des poids croissants \unsure{ds}.

\margin{\fbox{\parbox{.27\linewidth}{les poids sont-ils $\ge j$}} \deleted{\unsure{donc}}}

Si $X$ est \uline{propre}, on trouve que $H^i(X,F)$ est pur de poids $i+j$ (et semi-simple ici\ldots)!
